% ═══════════════════════════════════════════════════════════════
% 第三章：理论分析与研究假说
% ═══════════════════════════════════════════════════════════════

\section{理论分析与研究假说}

\subsection{理论框架}

本文的理论框架建立在信息不对称理论（Akerlof, 1970）、信号传递理论（Spence, 1973）、信息披露理论（Healy \& Palepu, 2001）和软信息价值理论（Liberti \& Petersen, 2019）的交叉基础之上。

在理想的有效市场中，上市公司股价反映所有公开信息，投资者可以通过价格信号做出合理决策。然而，现实中的资本市场存在严重的信息不对称。一方面，上市公司管理层比外部投资者掌握更多关于公司真实经营状况的信息，可能利用信息优势隐藏负面消息或进行盈余操纵。另一方面，上市公司的信息披露行为本身受到多种因素的影响，包括管理层动机、公司治理结构、外部监管压力和市场竞争环境等。

监管问询是监管者发现并纠正信息不对称问题的重要工具。当监管者（交易所）审查上市公司披露的公告和财务报告时，如果发现某些信号异常，便会发出问询函要求公司补充说明。从信息经济学的视角理解，监管问询行为本身是一个信息处理函数：

\begin{equation}
P(\text{Inquiry}_{i,t+{\Delta t}} = 1 \mid \Omega_{i,t}) = f(\text{Hard}_{i,t}, \text{Soft}_{i,t}, \text{Market}_{i,t}, \text{Governance}_{i,t}, \text{History}_{i,t})
\label{eq:framework}
\end{equation}

其中，$\Omega_{i,t}$是公司在时间$t$时点的全部可得信息集。$\text{Hard}_{i,t}$代表结构化财务指标（硬信息），$\text{Soft}_{i,t}$代表从公告文本中提取的语义特征（软信息）。$\text{Market}_{i,t}$和$\text{Governance}_{i,t}$代表市场和治理层面的控制变量，$\text{History}_{i,t}$代表历史监管记录。

本文的核心理论框架如图\ref{fig:theory_framework}所示。该框架区分了四类信息源在影响监管问询决策中的不同角色：硬信息直接反映财务健康状况的异常，软信息揭示管理层的潜在隐瞒动机和信息披露质量的隐性信号，市场信息反映外部参与者对公司风险的定价，治理信息反映内部监督机制的完备性。

\begin{figure}[H]
\centering
\small
\begin{tabular}{|c|}
\hline
\textbf{公司特征信息集} \\
\hline
\begin{tabular}{@{}c@{}}
\ \ 硬信息（财务指标）\ \ \\
\ \ 软信息（文本语义）\ \ \\
\ \ 市场信号（交易特征）\ \ \\
\ \ 治理信号（公司治理）\ \ \\
\hline
\end{tabular} \\
\downarrow \\
\begin{tabular}{@{}c@{}}
\textbf{信息不对称} \\
\textit{代理冲突 × 披露质量 × 外部监督} \\
\hline
\end{tabular} \\
\downarrow \\
\begin{tabular}{@{}c@{}}
\textbf{监管者决策函数} \\
\textit{信号发现 → 风险评估 → 问询决定} \\
\hline
\end{tabular} \\
\downarrow \\
\begin{tabular}{@{}c@{}}
\textbf{被解释变量：被问询概率} \\
$P(\text{Inquiry}_{i,t+60} = 1)$ \\
\hline
\end{tabular} \\
\hline
\end{tabular}
\caption{理论分析框架}
\label{fig:theory_framework}
\end{figure}

\subsection{研究假说}

\textbf{假说1（硬信息假说）：财务异常指标与监管问询概率显著正相关。}

传统的财务指标分析框架（如Beneish M-Score、Altman Z-Score）已被广泛用于识别公司风险和盈余操纵。这些指标直接反映了公司财务状况的异常程度。从监管者视角看，当公司的财务指标出现显著偏离行业均值或历史趋势的异常时，可能意味着公司存在收入确认问题、利润质量下降或潜在的财务舞弊行为，这会触发监管者的深入审查。具体而言，Beneish M-Score超过警戒线表明公司具有较强的盈余操纵动机，Altman Z-Score落入危险区表明公司面临较高的财务困境风险，应收增速显著超过收入增速则可能意味着收入质量下降。因此，我们预期这些财务异常指标与监管问询概率之间存在显著的正相关关系。

\textbf{假说2（软信息假说）：LLM提取的文本语义风险特征对监管问询概率具有显著的增量预测能力，且独立于财务指标。}

传统的财务指标（硬信息）具有可量化、可验证、标准化的优点，但同时也面临"硬信息陷阱"的局限——当财务报表经过合法合规编制时，单一的财务指标异常可能不足以触发问询。相比之下，公告文本中的语义特征（软信息）包含了管理层语调、信息披露充分性、MD\&A前瞻性陈述等丰富维度，这些维度难以通过简单的财务指标量化，但可能蕴含了关于公司真实风险状况的重要信号。大语言模型凭借其强大的语义理解能力，能够从非结构化文本中提取这些软信息特征，形成对硬信息的有效补充。

\textbf{假说3（异质性假说）：软信息的增量预测能力在信息不对称更严重的公司中更强。}

如果软信息的增量价值确实来源于其能够捕捉硬信息无法反映的信号，那么理论上，在信息不对称更严重的公司中，软信息的增量价值应当更大。这是因为，在信息不对称严重的公司中，硬信息的"信息充分性"更低，软信息作为补充渠道的价值相对更大。借鉴已有文献，我们选取以下两个维度度量信息不对称程度：公司市值规模（小市值公司信息不对称更严重）和分析师覆盖数量（低覆盖公司信息不对称更严重）。

\textbf{假说4（时间维度假说）：历史监管问询经历与当前问询概率显著正相关，监管问询行为具有时间序列自相关性。}

从监管实践看，历史被问询的公司往往在信息披露方面存在系统性缺陷，这类公司再次被问询的概率更高。同时，监管者对被问询公司的后续整改情况会保持持续关注，形成了"问询—整改—复查"的监管闭环。此外，同一行业的问询频率也会影响监管者的关注分配——当一个行业有多家公司近期被问询时，该行业的其他公司进入监管视野的概率也会提升。

\textbf{假说5（可解释性假说）：基于SHAP归因和LLM生成的预警报告可以提供有效的经济学归因，其解释质量满足业务复核要求。}

可解释性是本文的核心关注之一。我们将验证，结合SHAP数值归因与LLM自然语言生成的可解释报告机制是否能够产生逻辑完整、证据充分、业务可用的风险预警说明。借鉴LLM-as-Judge的评估框架\citep{zheng2024judge}，我们从风险识别准确性、证据充分性、逻辑链完整性、案例匹配合理性和业务实用性五个维度进行评估，预期综合得分不低于85分。